% Source: https://artofproblemsolving.com/community/c3h2665039_prove_there_exist_no_such_nonzero_integers
% Stable identity: aops-topic-2665039-q01
\textbf{(a)} Take $p=-2a$ and $q=-2b$, which are nonzero. Then $(a+p)^2-(b+q)^2=(-a)^2-(-b)^2=a^2-b^2=c^2$.

\textbf{(b)} Put $A=a+p$ and $B=b+q$, so we need $(A-B)(A+B)=c^2$ with $A+B>0$, hence $A-B>0$.

For $(5,4,3)$ with $p,q\ge1$: $A\ge6$ and $B\ge5$, so $A+B\ge11$. But $A-B\ge1$ forces $A+B\le9$. Contradiction, so no positive $p,q$ exist.

For $(10,6,8)$: $(A-B)(A+B)=64$ with $A+B\ge11+7=18$, and $A-B$, $A+B$ have the same parity. The factorisations of $64$ into two factors of equal parity with the larger at least $18$ are only $2\cdot32$ (since $4\cdot16$ has $16<18$). So $A=17$, $B=15$, and $\boxed{(p,q)=(7,9)}$ is the only pair: $17^2-15^2=289-225=64=8^2$.
